% from ega4-15-fr.tex
\oldpage[IV-3]{238}
se rencontrent pour $2\leqslant i\leqslant n$~; la réunion des $U_i\cap Z_{i-1}$ et des $U_i\cap Z_i$ pour $1\leqslant i\leqslant n$ est donc connexe et contient $x$ et $g(y)$. Comme $f|X'$ est universellement ouvert aux points de $X_y^{\prime 0}$, cela achève de prouver notre assertion, car si la proposition est prouvée lorsque $f$ est de type fini, la réunion $X^{\prime 0}$ des $X_z^{\prime 0}$ pour $z\in V$ sera un voisinage de $x$, et il en sera de même de $X^0$.

Supposons donc $f$ de type fini. Alors $X^0$ est localement constructible (9.7.12)~; il suffit par suite de montrer que $X^0$ contient toute générisation $x'$ de $x$ ($\mathbf{0}_{\mathrm{III}}$, 9.2.5). Posons $y'=f(x')$~; il existe un anneau de valuation discrète $A$ et un morphisme $u$ de $Y'=\operatorname{Spec}(A)$ dans $X$ tels que, si $z$ est le point fermé et $z'$ le point générique de $Y'$, on ait $u(z)=x$ et $u(z')=x'$ (\textbf{II}, 7.1.9)~; si $h=f\circ u$, on a donc $h(z)=y$ et $h(z')=y'$. Posons $X'=X\times_Y Y'$, $f'=f_{(Y')}:X'\to Y'$, qui est un morphisme de type fini universellement ouvert aux points de $f'^{-1}(z)$, et $g'=g_{(Y')}$, qui est une $Y'$-section de $X'$. Si $p_y:f'^{-1}(z)\to f^{-1}(y)$ (resp. $p_{y'}:f'^{-1}(z')\to f^{-1}(y')$) est la projection canonique, $p_y^{-1}(X_y^0)$ (resp. $p_{y'}^{-1}(X_{y'}^0)$) est la composante connexe de $f'^{-1}(z)$ (resp. $f'^{-1}(z')$) contenant $g'(z)$ (resp. $g'(z')$), car les $X_t^0$ sont géométriquement connexes (4.5.13) et il suffit d'appliquer (4.4.1). De plus, le morphisme $u$ correspond à une $Y'$-section $v$ de $X'$ telle que $p_y(v(y))=x$, $p_{y'}(v(y'))=x'$, de sorte que $v(y')$ est une générisation de $v(y)$, et il suffira donc de prouver que $v(y')$ appartient à $p_{y'}^{-1}(X_{y'}^0)$. Enfin, il est clair que $p_y^{-1}(X_y^0)$ est géométriquement réduit sur $\mathbf{k}(z)$.

En d'autres termes, on est ramené au cas où $Y$ est spectre d'un anneau de valuation discrète, $y$ son point fermé, $y'$ son point générique. Comme $f^{-1}(y)-X_y^0$ est alors fermé dans $X$, on peut, en remplaçant $X$ par l'ouvert $X-(f^{-1}(y)-X_y^0)$, supposer que $X_y^0=f^{-1}(y)$~; de même, on peut remplacer $X$ par l'ouvert, composante connexe de $X$, qui contient $g(Y)$, cette composante connexe contenant évidemment $X^0$~; autrement dit, on peut supposer $X$ connexe. La proposition sera alors établie si l'on montre que $X^0=X$, ou encore que $X_{y'}$ est connexe~; mais $f^{-1}(y')$ est géométriquement réduit sur $\mathbf{k}(y')$, et \emph{a fortiori} réduit~; comme en outre $f$ est ouvert aux points de $f^{-1}(y)$, on peut appliquer le lemme (15.5.6), d'où la conclusion.
\end{proof}

\begin{corollary}[15.6.5]
\label{IV.15.6.5-fr}
Soient $f:X\to Y$ un morphisme plat de présentation finie, $g$ une $Y$-section de $X$~; pour tout $y\in Y$, soit $X_y^0$ la composante connexe de $f^{-1}(y)$ contenant $g(y)$, et soit $X^0$ la réunion des $X_y^0$ quand $y$ parcourt $Y$. Alors, pour tout $y\in Y$ tel que $X_y^0$ soit géométriquement réduit sur $\mathbf{k}(y)$, $X^0$ est un voisinage de $X_y^0$ dans $X$. En particulier, si $f$ est réduit (6.8.1), $X^0$ est ouvert dans $X$.
\end{corollary}

\begin{proof}
La question étant locale sur $Y$, on peut supposer $Y=\operatorname{Spec}(A)$ affine. Il existe donc un sous-anneau noethérien $A_0$ et un morphisme plat de type fini $f_0:X_0\to Y_0=\operatorname{Spec}(A_0)$ tels que $X=X_0\times_{Y_0}Y$ et $f=(f_0)_{(Y)}$ (8.9.1 et 11.2.7)~; en outre, on peut supposer qu'il existe une $Y_0$-section $g_0$ de $X_0$ telle que $g=(g_0)_{(Y)}$ (8.8.2). Pour tout $y\in Y$, soit $y_0$ sa projection dans $Y_0$~; il résulte de (4.5.13) que $(X_0)_{y_0}^0$ est géométriquement connexe, donc, si $p:f^{-1}(y)\to f_0^{-1}(y_0)$ est la projection canonique, on a $X_y^0=p^{-1}((X_0)_{y_0}^0)$ (4.4.1)~; de plus, l'hypothèse que $X_y^0$ est géométriquement réduit sur $\mathbf{k}(y)$ entraîne que $(X_0)_{y_0}^0$ est géométriquement réduit sur $\mathbf{k}(y_0)$ (4.6.10). On est ainsi ramené au cas où on suppose

